\chapter{File Uploads (Multer + Cloudinary)}

\section{The Upload Pipeline}

Uploading an avatar or attachment has three parts: getting raw bytes off the
request (Multer), storing them on a CDN (Cloudinary), and saving the
resulting URL on a MongoDB document.

\begin{diagrambox}[Upload Pipeline]
\begin{verbatim}
React (FormData) --> POST multipart/form-data
        |
   Multer (parses body, exposes req.file)
        |
   Cloudinary (uploader.upload) --> secure_url
        |
   Save url on User/Task document
\end{verbatim}
\end{diagrambox}

\subsection{Why multipart/form-data}

JSON can only carry text. Binary files use \texttt{multipart/form-data}
encoding, which \texttt{express.json()} cannot parse --- hence a dedicated
middleware (Multer) for file routes.

\section{Backend: Multer Configuration}

Memory storage avoids an unnecessary disk write since the buffer is
forwarded straight to Cloudinary.

\begin{lstlisting}[language=JavaScript]
// backend/src/middleware/upload.middleware.js
const multer = require("multer");

const storage = multer.memoryStorage();

const fileFilter = (req, file, cb) => {
  const allowedMimeTypes = ["image/jpeg", "image/png", "image/webp"];
  if (allowedMimeTypes.includes(file.mimetype)) {
    cb(null, true);
  } else {
    cb(new Error("Only JPEG, PNG, or WEBP images are allowed"), false);
  }
};

const upload = multer({
  storage,
  fileFilter,
  limits: { fileSize: 2 * 1024 * 1024 }, // 2 MB max
});

module.exports = upload;
\end{lstlisting}

\subsection{Cloudinary Configuration}

\begin{lstlisting}[language=JavaScript]
// backend/src/config/cloudinary.js
const cloudinary = require("cloudinary").v2;

cloudinary.config({
  cloud_name: process.env.CLOUDINARY_CLOUD_NAME,
  api_key: process.env.CLOUDINARY_API_KEY,
  api_secret: process.env.CLOUDINARY_API_SECRET,
});

module.exports = cloudinary;
\end{lstlisting}

\begin{notebox}[NOTE]
Cloudinary credentials belong in \texttt{.env}, never committed to source
control (Chapter~24).
\end{notebox}

\subsection{Controller: Upload and Save}

\begin{lstlisting}[language=JavaScript]
// backend/src/controllers/user.controller.js
const cloudinary = require("../config/cloudinary");
const User = require("../models/User");

const uploadAvatar = async (req, res, next) => {
  try {
    if (!req.file) {
      return res.status(400).json({ message: "No image file provided" });
    }

    const base64 = req.file.buffer.toString("base64");
    const dataUri = `data:${req.file.mimetype};base64,${base64}`;

    const result = await cloudinary.uploader.upload(dataUri, {
      folder: "task-app/avatars",
      resource_type: "image",
    });

    const user = await User.findByIdAndUpdate(
      req.user.id,
      { avatar: result.secure_url },
      { new: true }
    ).select("-password");

    res.status(200).json({ user });
  } catch (error) {
    next(error);
  }
};

module.exports = { uploadAvatar };
\end{lstlisting}

\subsection{Route Wiring}

\begin{lstlisting}[language=JavaScript]
// backend/src/routes/user.routes.js
const express = require("express");
const upload = require("../middleware/upload.middleware");
const { uploadAvatar } = require("../controllers/user.controller");
const { protect } = require("../middleware/auth.middleware");

const router = express.Router();

router.post("/me/avatar", protect, upload.single("avatar"), uploadAvatar);

module.exports = router;
\end{lstlisting}

\begin{warnbox}[WARNING]
The field name in \texttt{upload.single("avatar")} must match the
\texttt{FormData} key exactly, or \texttt{req.file} is silently
\texttt{undefined}.
\end{warnbox}

\section{Frontend: Building and Sending FormData}

\begin{lstlisting}[language=JavaScript]
// frontend/src/components/AvatarUploader.jsx
import { useState } from "react";
import api from "../services/api";

function AvatarUploader({ onUploaded }) {
  const [file, setFile] = useState(null);
  const [uploading, setUploading] = useState(false);

  const handleSubmit = async (e) => {
    e.preventDefault();
    if (!file) return;

    const formData = new FormData();
    formData.append("avatar", file); // key MUST match upload.single("avatar")

    setUploading(true);
    try {
      const { data } = await api.post("/users/me/avatar", formData, {
        headers: { "Content-Type": "multipart/form-data" },
      });
      onUploaded(data.user.avatar);
    } catch (error) {
      console.error("Upload failed:", error.response?.data?.message);
    } finally {
      setUploading(false);
    }
  };

  return (
    <form onSubmit={handleSubmit}>
      <input
        type="file"
        accept="image/png, image/jpeg, image/webp"
        onChange={(e) => setFile(e.target.files[0])}
      />
      <button type="submit" disabled={!file || uploading}>
        {uploading ? "Uploading..." : "Upload Avatar"}
      </button>
    </form>
  );
}

export default AvatarUploader;
\end{lstlisting}

\begin{tipbox}[TIP]
Axios sets the \texttt{multipart/form-data} boundary automatically for a
\texttt{FormData} body; the explicit header is optional but harmless.
\end{tipbox}
